\section{Conclusion}
\label{sec:conclusion}

In this paper, we investigated whether the sparse, historically-limited training data of \talkie provides a more scientifically rigorous testbed for generalization than modern models. We hypothesized that the lack of contamination and the resulting ``Temporal Knowledge Cliff'' would create a clearer signal for structural generalization. Our experiments on \humaneval and \mmlu confirm this hypothesis: \talkie's failure in \zeroshot and success in \fewshot on modern tasks provide a binary-like indicator of reasoning that is often obscured in modern, dense-data models.

Our findings suggest that \talkie serves as an effective ``control group'' for the history of human knowledge. By eliminating the possibility of memorization for post-1930 concepts, we can isolate and measure the mechanics of machine logic in its pure form. This approach opens new avenues for studying how linguistic patterns from one era map to the technical requirements of another.

{\bf Future Work.} Future research should explore ``Analogical Transfer'' more deeply, for instance, by asking \talkie to explain modern quantum mechanics using only Newtonian analogies. Additionally, expanding the evaluation to more complex logic tasks beyond simple arithmetic would provide a better understanding of the scaling laws of generalization in historically-cutoff models.
\clearpage
